\documentclass[11pt,reqno]{amsart}

% --- Packages -----------------------------------------------------------------
\usepackage[T1]{fontenc}
\usepackage[utf8]{inputenc}
\usepackage[english]{babel}
\usepackage{amsmath,amssymb,amsthm,mathtools}
\usepackage{microtype}
\usepackage{booktabs}
\usepackage{enumitem}
\usepackage[dvipsnames]{xcolor}
\usepackage{hyperref}
\usepackage{xurl}
\hypersetup{
  colorlinks=true,
  linkcolor=blue!50!black,
  citecolor=blue!50!black,
  urlcolor=blue!50!black,
}

\raggedbottom

% --- Theorem environments -----------------------------------------------------
\theoremstyle{plain}
\newtheorem{theorem}{Theorem}[section]
\newtheorem{lemma}[theorem]{Lemma}
\newtheorem{corollary}[theorem]{Corollary}
\newtheorem{proposition}[theorem]{Proposition}
\newtheorem{conjecture}[theorem]{Conjecture}

\theoremstyle{definition}
\newtheorem{definition}[theorem]{Definition}

\theoremstyle{remark}
\newtheorem{remark}[theorem]{Remark}

% --- Macros -------------------------------------------------------------------
\newcommand{\Tab}{T_{a,b}}
\newcommand{\Tnegq}{T_{-q,b}}
\newcommand{\Tneg}{T_{-}}
\newcommand{\Tpos}{T_{+}}
\newcommand{\Z}{\mathbb{Z}}
\newcommand{\N}{\mathbb{N}}
\newcommand{\Q}{\mathbb{Q}}
\newcommand{\Qtwo}{\mathbb{Q}_{2}}
\newcommand{\Ztwo}{\mathbb{Z}_{2}}
\newcommand{\Nodd}{\N_{\mathrm{odd}}}
\newcommand{\Obs}{\mathcal{O}}
\newcommand{\Atom}{\mathcal{A}}
\newcommand{\Vk}{V_{K}}
\newcommand{\Vi}{V_{I}}
\newcommand{\vk}{v_{K}}
\newcommand{\vi}{v_{I}}
\newcommand{\aK}{a_{K}}
\newcommand{\aI}{a_{I}}
\newcommand{\vtwo}{v_{2}}
\newcommand{\vM}{v_{M}}
\newcommand{\ord}{\mathrm{ord}}
\newcommand{\cWq}{c_{W}^{(-q)}}
\newcommand{\Obsq}{\Obs_{L}^{(-q,b)}}
\newcommand{\Atomq}{\Atom_{L}^{(-q,b)}}
\newcommand{\srcpath}[1]{\nolinkurl{#1}}
\emergencystretch=3em

% --- Document -----------------------------------------------------------------
\title[Obstruction residues at negative-Mersenne multipliers]{Obstruction
  residues at negative-Mersenne multipliers:\\
  a synchronisation hierarchy and density bounds}

\author{Jakob Stemberger}
\address{Independent researcher, Vienna, Austria}
\email{jakob.stemberger@yaccob.net}
\urladdr{https://orcid.org/0009-0001-2746-6877}
\date{\today}

\subjclass[2020]{11B83, 11S05, 37P05, 37B10, 37B40}

\keywords{Collatz conjecture, $3n\pm1$ map, $(an+b)$ family, Mersenne
  multiplier, residue dynamics, parallel reduction, 2-adic density}

\begin{document}

\begin{abstract}
For a Mersenne multiplier $q = 2^{\vM} - 1$ with $\vM \ge 2$ and an odd
bias $b \ne 0$, the accelerated map $\Tnegq(n) := (-qn + b)/2^{\vtwo(-qn+b)}$
acts on the odd integers, and the companion $(an+b)$-family manuscript
singles it out as the one case admitting a finer structure than the general
theory captures. We develop that structure. The starting point is an exact
algebraic identity, $-q r + b = 2^v(-q m + b) + (2^{\vM} - 2^v) b$ whenever
$r + b = 2^v m$, whose correction term $(2^{\vM} - 2^v) b$ vanishes precisely
at $v = \vM$; the vanishing is exclusive to Mersenne multipliers and
synchronises a \emph{principal class} of obstructions. Iterating the
two-track parallel reduction is governed by a linear recursion on an affine
certificate $\aK^{(k)} = \lambda_k \aI^{(k)} + D_k b$, whose synchronising
states are parametrised by a signed index $j \in \Z$ as
$(\lambda_k, D_k) = (2^{j\vM}, -(2^{j\vM}-1)/q)$. This yields a sub-class
hierarchy indexed by the synchronisation level $k$ and the signed index $j$:
we classify the level-$3$ sub-classes (three universal families for
$\vM \ge 3$, four at $\vM = 2$) with explicit substitution chains, exhibit
the level-$4$ families (including a genuinely signed $j = -1$ family with
a rational synchronisation), and prove a node-lift theorem showing that
inserting a fixed-point step at the unique fixed point
$(\lambda^*, D^*) = (2^{\vM+1}, -1)$ of the recursion lifts each sub-class to
the next level, generating four iterated families. Summing the resulting
disjoint residue cylinders gives the density lower bound
$\cWq \ge 1/(2q)$, which we refine; the same families yield a matching upper
bound, making $1/(2q)$ asymptotically sharp as $\vM \to \infty$ ---
$\cWq = (1/(2q))(1 + O(2^{-\vM}))$ --- conditional on a
classification-completeness hypothesis. A companion
observability bound $\lvert j \rvert \le (L - k)/\vM$ controls the level at
which a sub-class first becomes visible. Each finitely-checkable claim is
backed by an executable falsifier --- a check that surfaces a counterexample
if the claim is false; the limit and asymptotic statements are
recorded as such.
\end{abstract}

\maketitle

% =============================================================================
\section{Introduction}\label{sec:intro}
% =============================================================================

\subsection{Setting}

For odd integers $a, b$ with $\lvert a \rvert \ge 3$ and $b \ne 0$, the
\emph{affine accelerated map} $\Tab$ acts on the odd integers by
\[
\Tab(n) := \frac{an + b}{2^{\vtwo(an + b)}}.
\]
For $a > 0$ this restricts to a self-map of the positive odd integers
$\Nodd$; for $a < 0$ --- the case of this paper --- an orbit may leave the
positives, but the obstruction analysis is purely residue-theoretic (the
obstruction sets live in $\{1, \ldots, 2^L\}$, and Section~\ref{sec:density}
reads their density as a Haar measure on the $2$-adic integers
$\Ztwo$), so the sign of intermediate values is immaterial.
Two specialisations are classical: $(a,b) = (3,1)$ is the Collatz map
$\Tpos$ \cite{Lagarias2010,Lagarias2010bib}, and $(a,b) = (3,-1)$ is the
$3n-1$ map $\Tneg$ \cite{Lagarias2010bib}. Generalised $(an+b)$ maps go back
to Hasse's generalised Syracuse algorithm, whose theory is developed by
Matthews and Watts \cite{MatthewsWatts1984} and M\"oller \cite{Moeller1978};
the broader class of piecewise-defined generalised Collatz iterations is
computationally universal \cite{Conway1972}, which is part of what makes an
exact, level-by-level algebraic characterisation of a structured subfamily
worthwhile. The strongest unconditional progress towards the Collatz
conjecture is Tao's theorem that almost all orbits attain almost bounded
values \cite{Tao2022}.

This paper studies the \emph{negative-Mersenne multipliers}
\[
a = -q, \qquad q := 2^{\vM} - 1 \ \text{Mersenne}, \quad \vM \ge 2,
\]
so $\lvert a \rvert \in \{3, 7, 15, 31, 63, \ldots\}$. The companion
manuscript on the $(an+b)$ family \cite{ObstructionResiduesFamily2026}
develops a uniform theory of obstruction residues across all odd multipliers
with $\lvert a \rvert \ge 3$ --- a bias-axis conjugation that makes
cardinality and density depend only on $a$, a multiplier-order lemma, the
existence of the asymptotic density, and a no-shift-zero-atom theorem --- and
explicitly flags the negative-Mersenne case as admitting a \emph{finer}
sub-class hierarchy, treated ``in a separate manuscript.'' The present paper
is that manuscript. The base case $(a,b) = (3,-1)$, i.e.\ $\vM = 2$, is the
$\Tneg$-manuscript \cite{ObstructionResidues2026}, which proves a constructive
lift, positive density, and full factor complexity; here we explain what is
special about the whole Mersenne ladder $\vM = 2, 3, 4, \ldots$.

\subsection{Obstruction residues}

We recall the obstruction-residue framework of
\cite{ObstructionResiduesFamily2026} in the specialisation $a = -q$; a
self-contained account is given in Section~\ref{sec:setup}. For an odd
residue $r$ and a level $L$, a two-track parallel reduction runs the map
$\Tnegq$ simultaneously on $r$ and on its first reduced value $m$, carrying
an affine certificate between the tracks; $r$ is an \emph{obstruction
residue} at level $L$ when the certificate vanishes at termination, i.e.\
when the two tracks synchronise. The set of such residues in
$\{1, \ldots, 2^L\}$ is denoted $\Obsq$, the \emph{atomic} obstructions
(those not lifted from level $L-1$) are denoted $\Atomq$, and the asymptotic
density is
\[
\cWq := \lim_{L \to \infty} \frac{\lvert \Obsq \rvert}{2^L},
\]
which exists by \cite{ObstructionResiduesFamily2026}. The fine invariant of
an obstruction is its \emph{synchronisation level} $k$ --- the index at which
the two tracks first agree --- together with a \emph{signed synchronisation
index} $j \in \Z$ recording the cumulative valuation gap between the tracks.

\subsection{Main results}

The Mersenne condition enters through a single exact identity
(Theorem~\ref{thm:main-identity}): for $r + b = 2^v m$,
\[
-q r + b = 2^v(-q m + b) + (2^{\vM} - 2^v)\, b,
\]
whose correction term $(2^{\vM} - 2^v) b$ vanishes \emph{exactly} at
$v = \vM$, and only for Mersenne $q$ does such a $v$ exist for arbitrary $b$
(Corollaries~\ref{cor:mersenne-sync} and~\ref{cor:nonmersenne}). The
vanishing produces the \emph{principal synchronising class} --- the residues
with $\vtwo(r + b) = \vM$, which synchronise in a single step; iterating it is
governed by a linear recursion on the affine certificate
(Theorem~\ref{thm:recursion}), whose synchronising states are exactly
$(\lambda_k, D_k) = (2^{j\vM}, -(2^{j\vM}-1)/q)$, $j \in \Z$.

This recursion structures the obstructions into a sub-class hierarchy:
\begin{itemize}[leftmargin=1.4em]
\item a level-$3$ structure theorem (Theorem~\ref{thm:stage3}): for
  $\vM \ge 3$ exactly three universal families $A$, $B$, $C$, with explicit
  substitution chains, and a fourth family at $\vM = 2$;
\item the level-$4$ families (Theorem~\ref{thm:stage4}), including a
  genuinely signed family $G_1$ with $j = -1$ that synchronises rationally;
\item a node-lift theorem (Theorem~\ref{thm:node-lift}): the recursion has a
  unique fixed point $(\lambda^*, D^*) = (2^{\vM+1}, -1)$, and inserting the
  fixed-point step at a node lifts a sub-class to the next level, generating
  four iterated families $A_k, B_k, C_k, F_{1,k}$;
\item a density lower bound $\cWq \ge 1/(2q)$
  (Theorem~\ref{thm:atom-density}), obtained by summing disjoint residue
  cylinders, refined in Corollary~\ref{cor:atom-density-extended} and shown
  asymptotically sharp under a classification-completeness hypothesis
  (Theorem~\ref{thm:asymptotics}): $\cWq = (1/(2q))(1 + O(2^{-\vM}))$;
\item an observability bound $\lvert j \rvert \le (L - k)/\vM$
  (Theorem~\ref{thm:j-bound}) controlling the level at which a level-$k$,
  index-$j$ sub-class first becomes visible.
\end{itemize}

The density scale deserves emphasis. With $\lvert a \rvert = 2^{\vM} - 1$ the
multiplicative order $e := \ord_{\lvert a \rvert}(2)$ equals $\vM$, so the
bound $1/(2q) = \Theta(2^{-\vM}) = \Theta(2^{-e})$ is the \emph{square root}
of the empirical $\Theta(4^{-e})$ density observed at generic positive
multipliers in \cite{ObstructionResiduesFamily2026}: negative-Mersenne
multipliers are anomalously obstruction-rich, by an exponential factor, and
the present hierarchy is the structural reason.

\subsection{Relation to the family manuscript}

We use the obstruction-residue framework, the parallel reduction, and the
existence of the asymptotic density from
\cite{ObstructionResiduesFamily2026} as a black box (recalled in
Section~\ref{sec:setup}); everything specific to the Mersenne ladder ---
the main identity, the recursion, the sub-class families, the node lift, and
the density bounds --- is proved here. The family manuscript's
no-shift-zero-atom theorem holds ``outside the negative-Mersenne case'': its
proof forces a shift-zero obstruction to take $J \ge 2$ steps unless
$2^v = \lvert a \rvert + 1$ is solvable, which happens exactly for Mersenne
$\lvert a \rvert$ (at $v = \vM$). That one-step exception is the principal
class of the present paper, so the family theorem excludes precisely the
negative-Mersenne ladder and the two accounts are complementary rather than
overlapping.

\subsection{Outline}

Section~\ref{sec:setup} fixes the parallel reduction and the affine
certificate. Section~\ref{sec:identity} proves the main identity and its
Mersenne/non-Mersenne corollaries. Section~\ref{sec:recursion} derives the
recursion theorem and identifies its synchronising states.
Sections~\ref{sec:stage3} and~\ref{sec:stage4} classify the level-$3$ and
level-$4$ sub-classes. Section~\ref{sec:nodelift} proves the node-lift
theorem. Section~\ref{sec:density} establishes the density lower bound and
its refinement, Section~\ref{sec:jbound} the observability bound, and
Section~\ref{sec:asymptotics} the asymptotic sharpness.
Section~\ref{sec:open} collects open problems.

% =============================================================================
\section{Setup: parallel reduction and the affine certificate}\label{sec:setup}
% =============================================================================

Fix throughout a Mersenne multiplier $q = 2^{\vM} - 1$ with $\vM \ge 2$ and an
odd bias $b \ne 0$; we write $a = -q$ and study $\Tnegq$. The
family-existence results and the density lower bounds below hold for any such
$b$; the \emph{converse} classification of synchronising states
(Theorem~\ref{thm:recursion}) and the exact sub-class counts it yields
(Theorem~\ref{thm:stage3}) assume in addition that $b$ is coprime to $q$ ---
satisfied by the classical biases $b = \pm 1$ and by every bias used in the
verification.

\subsection{Nesting and the two-track reduction}

Let $r$ be a positive odd integer. Since $b$ is odd, $r + b$ is even; write
\[
v_0 := \vtwo(r + b) \ge 1, \qquad m_0 := (r + b)/2^{v_0} \ \text{(odd)}.
\]
The \emph{parallel reduction} at level $L > v_0$ runs two tracks. The
\emph{K-track} carries the value $\aK^{(0)} := r$ and modulus exponent $L$;
the \emph{I-track} carries $\aI^{(0)} := m_0$ and modulus exponent
$L - v_0$. At each step $k \ge 0$ set
\[
\vk^{(k)} := \vtwo\!\big(-q\,\aK^{(k)} + b\big), \qquad
\vi^{(k)} := \vtwo\!\big(-q\,\aI^{(k)} + b\big),
\]
and update
\[
\aK^{(k+1)} := \frac{-q\,\aK^{(k)} + b}{2^{\vk^{(k)}}}, \qquad
\aI^{(k+1)} := \frac{-q\,\aI^{(k)} + b}{2^{\vi^{(k)}}},
\]
each step being admissible while the corresponding modulus exponent stays
positive. The cumulative valuations are
$\Vk^{(k)} := \sum_{i < k} \vk^{(i)}$ and
$\Vi^{(k)} := \sum_{i < k} \vi^{(i)}$, so the residual modulus exponents are
$L - \Vk^{(k)}$ and $L - v_0 - \Vi^{(k)}$.

\subsection{The affine certificate}

\begin{lemma}[Affine certificate]\label{lem:certificate}
For every step $k \ge 0$ there are $\lambda_k \in \Q^{\times}$ and
$D_k \in \Q$ with
\[
\aK^{(k)} = \lambda_k\, \aI^{(k)} + D_k\, b \qquad \text{in } \Q,
\]
and $\lambda_0 = 2^{v_0}$, $D_0 = -1$.
\end{lemma}

\begin{proof}
From $r + b = 2^{v_0} m_0$ we get $r = 2^{v_0} m_0 - b$, i.e.\
$\aK^{(0)} = 2^{v_0}\aI^{(0)} - b$, the claim at $k = 0$ with
$(\lambda_0, D_0) = (2^{v_0}, -1)$. For the inductive step, substitute the
affine form $\aK^{(k)} = \lambda_k \aI^{(k)} + D_k b$ into
$-q\,\aK^{(k)} + b$ and divide by $2^{\vk^{(k)}}$: this rewrites
$\aK^{(k+1)}$ once more as a $\Q$-affine function of $\aI^{(k+1)}$ and $b$,
so the form persists with some $(\lambda_{k+1}, D_{k+1}) \in
\Q^{\times} \times \Q$. The explicit update for $(\lambda_{k+1}, D_{k+1})$
is computed in Theorem~\ref{thm:recursion}.
\end{proof}

\begin{definition}[Synchronisation level and signed index]\label{def:sync}
The reduction \emph{synchronises} at the least index $k$ with
$\aK^{(k)} = \aI^{(k)}$, equivalently $(\lambda_k, D_k) = (1, 0)$; this $k$ is
the \emph{synchronisation level}. The \emph{signed synchronisation index} is
the unique integer $j$ for which the state one step before synchronisation
has $\lambda_{k-1} = 2^{j\vM}$; this is well-defined because
Theorem~\ref{thm:recursion} shows every pre-synchronisation $\lambda$ is the
integral power $2^{j\vM}$, $j \in \Z$. A residue $r$ is an
\emph{obstruction residue} at level $L$, $r \in \Obsq$, when the two tracks
synchronise before either modulus exponent is exhausted; $\Atomq \subseteq
\Obsq$ are those not arising as a lift of a level-$(L-1)$ obstruction, and
\[
\cWq := \lim_{L \to \infty} \lvert \Obsq \rvert / 2^L .
\]
\end{definition}

\begin{remark}
Definition~\ref{def:sync} specialises the obstruction residues of
\cite{ObstructionResiduesFamily2026} to $a = -q$: there the certificate is
written $\aI = c\,\aK + d$ and the vanishing condition $\Phi_J = 0$ is the
synchronisation $\aK = \aI$ here. The family's shift index $\delta$ (the
valuation gap at termination) and the signed synchronisation index $j$ used
here (read off one step before synchronisation) are different invariants.
What the family's no-shift-zero-atom theorem excludes \emph{outside} the
negative-Mersenne case is a one-step ($J = 1$) shift-zero obstruction: its
proof forces $J \ge 2$ unless $2^v = \lvert a \rvert + 1$ is solvable, which
holds exactly for Mersenne $\lvert a \rvert$ at $v = \vM$. That solvable case
is the principal Mersenne-synchronising class here --- which carries
$j = 1$, not $j = 0$. The Mersenne ladder populates $j \in \Z$ (the level-3 families have
$j \in \{1, 2\}$, the level-4 family $G_1$ has $j = -1$), so the signedness
of $j$ is essential. Existence of the limit $\cWq$ is the family's
density-limit proposition, used as a black box.
\end{remark}

\begin{definition}[Universal vs.\ $q$-specific sub-classes]\label{def:universal}
A \emph{sub-class} is a set of obstruction residues sharing one initial
nesting exponent $v_0$ and one driving sequence of valuation pairs
$(\vk^{(i)}, \vi^{(i)})$. The sub-class is \emph{universal} when this
valuation data is \emph{affine in $\vM$} --- each of $v_0, \vk^{(i)},
\vi^{(i)}$ is an integer-affine function $\alpha\vM + \beta$ with fixed
integers $\alpha \ge 0$ and $\beta$ independent of $\vM$ (so a constant
$\beta$, the offset form $\vM + \beta$, or e.g.\ $2\vM$ are all admitted)
--- so that a single closed form describes it for all $\vM$ in its range;
otherwise it is \emph{$q$-specific}. The exhaustiveness statements below
(``exactly three families'', etc.) are understood \emph{within the universal
sub-classes}: they enumerate the affine-in-$\vM$ valuation patterns and do
not assert an a-priori bound ruling out non-affine families. The number of
$q$-specific sub-classes is itself $\vM$-dependent
(Theorem~\ref{thm:stage4}, Conjecture~\ref{conj:qspecific}).
\end{definition}

% =============================================================================
\section{The main identity}\label{sec:identity}
% =============================================================================

\begin{theorem}[Main identity]\label{thm:main-identity}
Let $q = 2^{\vM} - 1$. For all integers $r, m, b$ and all $v \ge 1$ with
$r + b = 2^v m$,
\[
-q r + b = 2^v\,(-q m + b) + \big(2^{\vM} - 2^v\big)\, b .
\]
\end{theorem}

\begin{proof}
From $r = 2^v m - b$,
\[
-q r + b = -q\big(2^v m - b\big) + b = -q\,2^v m + (q + 1)\, b .
\]
Since $q + 1 = 2^{\vM}$,
\[
-q\,2^v m + 2^{\vM} b = 2^v(-q m) + 2^v b + \big(2^{\vM} - 2^v\big) b
= 2^v(-q m + b) + \big(2^{\vM} - 2^v\big) b . \qedhere
\]
\end{proof}

\begin{corollary}[Mersenne synchronisation]\label{cor:mersenne-sync}
At $v = \vM$ the correction term vanishes and
$-q r + b = 2^{\vM}(-q m + b)$. Moreover, among all $q \ge 2$, a level $v$
with $(q + 1 - 2^v)\, b = 0$ for arbitrary $b$ exists if and only if
$q = 2^{\vM} - 1$ is Mersenne, in which case $v = \vM$ is the unique such
level.
\end{corollary}

\begin{proof}
Setting $v = \vM$ gives $2^{\vM} - 2^v = 0$, hence the displayed identity.
For the characterisation: $(q + 1 - 2^v) b = 0$ for $b \ne 0$ forces
$q + 1 = 2^v$, i.e.\ $q + 1$ is a power of two; this holds for some $v$ iff
$q$ is Mersenne, and then, by uniqueness of the binary representation,
$v = \vM = \vtwo(q+1)$ is the only such level.
\end{proof}

\begin{corollary}[Non-Mersenne asymmetry]\label{cor:nonmersenne}
For an arbitrary integer $q \ge 2$ and $r + b = 2^v m$,
\[
-q r + b = 2^v(-q m + b) + (q + 1 - 2^v)\, b .
\]
If $q$ is not Mersenne, then $q + 1 - 2^v \ne 0$ for every $v \ge 1$, so for
every odd $b \ne 0$ the correction term is nonzero at all levels and no
principal synchronising class exists.
\end{corollary}

\begin{proof}
The displayed identity is the computation in the proof of
Theorem~\ref{thm:main-identity} with $q + 1$ in place of $2^{\vM}$. If $q$ is
not Mersenne then $q + 1$ is not a power of two, so $q + 1 \ne 2^v$ for all
$v$; thus $(q + 1 - 2^v) b \ne 0$ for odd $b \ne 0$.
\end{proof}

The finite arithmetic of Theorem~\ref{thm:main-identity} and
Corollaries~\ref{cor:mersenne-sync}--\ref{cor:nonmersenne} is exercised over
many samples by \srcpath{code/python/identity.py}.

% =============================================================================
\section{The recursion theorem}\label{sec:recursion}
% =============================================================================

\begin{theorem}[Recursion theorem]\label{thm:recursion}
Let $q = 2^{\vM} - 1$, and suppose $\aK^{(k)} = \lambda_k \aI^{(k)} + D_k b$.
(The converse classification of synchronising states below assumes in
addition $\gcd(b, q) = 1$; cf.\ Section~\ref{sec:setup}.) Then
\[
\lambda_{k+1} = \lambda_k \cdot 2^{\,\vi^{(k)} - \vk^{(k)}}, \qquad
D_{k+1} = \frac{1 - \lambda_k - q\,D_k}{2^{\vk^{(k)}}} ,
\]
with initial data $(\lambda_0, D_0) = (2^{v_0}, -1)$. Integrality of
$D_{k+1}$ is the consistency condition
$2^{\vk^{(k)}} \mid E_k$ with $E_k := 1 - \lambda_k - q D_k$. The reduction
synchronises ($\aK^{(k)} = \aI^{(k)}$) iff $(\lambda_k, D_k) = (1, 0)$, and
the states one step before synchronisation are, assuming $\gcd(b, q) = 1$,
exactly
\[
(\lambda_k, D_k) = \Big(2^{\,j\vM},\ -\tfrac{2^{\,j\vM} - 1}{q}\Big),
\qquad j \in \Z,
\]
each of which is integral because $q = 2^{\vM} - 1 \mid 2^{j\vM} - 1$.
\end{theorem}

\begin{proof}
Substituting $\aK^{(k)} = \lambda_k \aI^{(k)} + D_k b$ into $-q\aK^{(k)} + b$
and using $q + 1 = 2^{\vM}$ only at the end,
\[
-q\,\aK^{(k)} + b
= \lambda_k\big(-q\,\aI^{(k)} + b\big) + \big(1 - \lambda_k - q D_k\big) b
= \lambda_k\big(-q\,\aI^{(k)} + b\big) + E_k\, b .
\]
By definition of the step, $-q\,\aK^{(k)} + b = 2^{\vk^{(k)}} \aK^{(k+1)}$
and $-q\,\aI^{(k)} + b = 2^{\vi^{(k)}} \aI^{(k+1)}$; hence
\[
2^{\vk^{(k)}} \aK^{(k+1)}
= \lambda_k\, 2^{\vi^{(k)}} \aI^{(k+1)} + E_k\, b ,
\]
and dividing by $2^{\vk^{(k)}}$ gives the stated updates with
$\lambda_{k+1} = \lambda_k 2^{\vi^{(k)} - \vk^{(k)}}$ and
$D_{k+1} = E_k / 2^{\vk^{(k)}}$. Since $\aK^{(k+1)}, \aI^{(k+1)} \in \Z$, the
quantity $D_{k+1}$ is an integer iff $2^{\vk^{(k)}} \mid E_k$. The initial
data are Lemma~\ref{lem:certificate}.

Synchronisation $\aK^{(k)} = \aI^{(k)}$ for all admissible $\aI^{(k)}$ is
$\lambda_k = 1$, $D_k = 0$. Now take a state $(\lambda_k, D_k) =
(2^{j\vM}, -(2^{j\vM} - 1)/q)$. Then
\[
E_k = 1 - 2^{j\vM} - q\Big(\!-\tfrac{2^{j\vM} - 1}{q}\Big)
   = 1 - 2^{j\vM} + (2^{j\vM} - 1) = 0 ,
\]
so $D_{k+1} = 0$; choosing the step with $\vi^{(k)} - \vk^{(k)} = -j\vM$
gives $\lambda_{k+1} = 2^{j\vM} \cdot 2^{-j\vM} = 1$, i.e.\ the next step
synchronises.

Conversely --- here we use $\gcd(b, q) = 1$ (the standing coprimality
assumption of Section~\ref{sec:setup}) --- suppose the next step
synchronises: $\lambda_{k+1} = 1$ and $D_{k+1} = 0$, with common synchronised
value $w := \aK^{(k+1)} = \aI^{(k+1)}$. From $\lambda_{k+1} = \lambda_k 2^{\vi^{(k)} - \vk^{(k)}} = 1$ we
get $\lambda_k = 2^{s}$ with $s := \vk^{(k)} - \vi^{(k)} \in \Z$, and
$E_k = 0$ gives $D_k = (1 - 2^{s})/q$. We show $\vM \mid s$ \emph{for either
sign of $s$}, so that $j := s/\vM \in \Z$. Inverting the synchronising step,
\[
\aK^{(k)} = \frac{b - 2^{\vk^{(k)}} w}{q}, \qquad
\aI^{(k)} = \frac{b - 2^{\vi^{(k)}} w}{q},
\]
are both integers, so $q$ divides their difference
$(2^{\vi^{(k)}} - 2^{\vk^{(k)}})\, w$. Writing $m := \min(\vk^{(k)},
\vi^{(k)})$, this difference equals $\pm\, 2^{m}(2^{\lvert s\rvert} - 1)\, w$,
and $w$ is itself coprime to $q$ --- a common prime factor of $w$ and $q$
would divide $q\,\aK^{(k)} = b - 2^{\vk^{(k)}} w$, hence $b$, contradicting
$\gcd(b, q) = 1$. Since $q$ is odd, it follows that $q \mid 2^{\lvert
s\rvert} - 1$. Now $2$ is a unit of multiplicative order exactly $\vM$
modulo $q$: indeed $2^{\vM} = q + 1 \equiv 1 \pmod q$, while
$0 < 2^t - 1 < 2^{\vM} - 1 = q$ for $0 < t < \vM$, so that
$2^t \equiv 1 \pmod q$ holds if and only if $\vM \mid t$, for every integer
$t \ge 0$. Applying this with $t = \lvert s\rvert$ gives $\vM \mid s$, whence
$\lambda_k = 2^{j\vM}$ and $D_k = -(2^{j\vM} - 1)/q$ with $j \in \Z$ --- the
negative-$j$ states being the rational ones ($\lambda_k = 2^{j\vM} < 1$)
realised by the $G_1$ family of Section~\ref{sec:stage4}.
\end{proof}

The recursion and its consistency arithmetic are checked over many step
sequences by \srcpath{code/python/recursion.py}.

% =============================================================================
\section{The level-3 structure theorem}\label{sec:stage3}
% =============================================================================

\begin{theorem}[Level-3 structure]\label{thm:stage3}
Let $q = 2^{\vM} - 1$ with $\gcd(b, q) = 1$ (Section~\ref{sec:setup}), so
that the converse of Theorem~\ref{thm:recursion} applies. For $\vM \ge 3$ the
obstruction residues that
synchronise at level $k = 3$ fall into exactly three universal families,
with the data
\[
\renewcommand{\arraystretch}{1.2}
\begin{array}{c|c|c|c|c}
\text{family} & v_0 & \big(\vk^{(0)},\vi^{(0)}\big),\,\big(\vk^{(1)},\vi^{(1)}\big)
 & j & (\lambda_2, D_2) \\ \hline
A & \vM + 1 & (\vM, \vM),\,(\vM, \vM - 1) & 1 & (2^{\vM}, -1) \\
B & \vM - 1 & (\vM - 1, 1),\,(\vM, 2\vM - 1) & 1 & (2^{\vM}, -1) \\
C & 1       & (1, \vM + 1),\,(\vM - 1, 2\vM - 2) & 2 & (2^{2\vM}, -(2^{\vM}+1))
\end{array}
\]
and synchronisation identity
$\aK^{(2)} - 2^{j\vM}\aI^{(2)} = -\frac{2^{j\vM} - 1}{q}\, b$
(equal to $-b$ for $A, B$ and $-(2^{\vM} + 1)\, b$ for $C$). The residues of
families $A$ and $B$ are the arithmetic progressions
\[
\begin{aligned}
A:\ &\ r = 2^{3\vM}\nu - K_A\, b, & K_A &= 2^{2\vM + 2} + 2^{\vM + 1} + 1, & \nu &\ \text{odd};\\
B:\ &\ r = 2^{3\vM}\tau - K_B\, b, & K_B &= 5\cdot 2^{2\vM - 1} + 3\cdot 2^{\vM - 1} + 1, & \tau &\in \Z.
\end{aligned}
\] The constraint $\nu$ odd in family $A$ selects its atomic residues --- those
whose leading nesting coefficient is odd, halving the cylinder and giving the
per-level density $2^{-(k\vM + 1)}$ --- whereas family $B$'s cylinder is full
($\tau \in \Z$, per-level density $2^{-k\vM}$); this asymmetry is the source of
the differing family contributions to the density bound of
Section~\ref{sec:density}. For $\vM = 2$ there is a fourth family
$D$ with $v_0 = \vM + 3 = 5$ and $j = 3$.
\end{theorem}

\begin{proof}
A level-$3$ sub-class (Definition~\ref{def:universal}) is a choice of initial
nesting exponent $v_0$ and a sequence of valuation pairs
$(\vk^{(0)}, \vi^{(0)}), (\vk^{(1)}, \vi^{(1)})$
driving the recursion of Theorem~\ref{thm:recursion} from $(\lambda_0, D_0) =
(2^{v_0}, -1)$ to a state $(\lambda_2, D_2)$ one step before
synchronisation, hence of the form $(2^{j\vM}, -(2^{j\vM} - 1)/q)$. Running
the recursion symbolically in $\vM$ for each of the three displayed
$(v_0, \text{pair-sequence})$ data yields exactly the listed $(\lambda_2,
D_2)$, and back-substituting through the nesting recovers the stated
arithmetic progressions; the synchronisation identity is the relation
$\aK^{(2)} = \lambda_2 \aI^{(2)} + D_2 b$ read off from $(\lambda_2, D_2)$.

For completeness we enumerate the admissible \emph{universal} valuation data
(Definition~\ref{def:universal}, affine in $\vM$; the families above realise
the offsets $0, \pm 1, \pm 2$). By
Theorem~\ref{thm:recursion} the two steps must drive $(2^{v_0}, -1)$ to a
pre-synchronisation state $(2^{j\vM}, -(2^{j\vM}-1)/q)$, and by the
observability bound (Theorem~\ref{thm:j-bound}) a level-$3$ sub-class visible
at any level has $\lvert j\rvert$ bounded; the consistency relations
$2^{\vk^{(i)}} \mid E_i$ then leave finitely many affine offset patterns to
test, each a polynomial identity in $\vM$. Enumerating them --- symbolically
in $\vM$ (which yields exactly $A, B, C$ for $\vM \ge 3$) and concretely for
$q = 3, 7$ via \srcpath{code/python/stage3.py} --- shows there are no others. At
$\vM = 2$ the families $B$ ($v_0 = \vM - 1 = 1$) and $C$ ($v_0 = 1$) share the
same initial nesting exponent --- the offset $\vM - 2$ separating their
nesting depths collapses to $0$ --- and an additional valuation pattern with
$v_0 = 5$ becomes admissible, giving the extra family $D$ ($j = 3$);
the concrete counts confirm this --- four sub-classes summing to $132$
obstruction residues at $L = 12$ for $q = 3$, three summing to $64$ at
$L = 14$ for $q = 7$.
The general-$\vM$ exhaustiveness rests on this affine-shape enumeration
rather than on an a-priori bound ruling out non-affine families.
\end{proof}

% =============================================================================
\section{The level-4 families}\label{sec:stage4}
% =============================================================================

\begin{theorem}[Level-4 families]\label{thm:stage4}
Let $q = 2^{\vM} - 1$ with $\vM \ge 3$. There are at least five universal
families synchronising at level $k = 4$:
\[
\renewcommand{\arraystretch}{1.2}
\begin{array}{c|c|c|c}
\text{family} & v_0 & \big(\vk^{(0)},\vi^{(0)}\big),\,\big(\vk^{(1)},\vi^{(1)}\big),\,\big(\vk^{(2)},\vi^{(2)}\big) & j \\ \hline
A_{\mathrm{lift}} & \vM + 1 & (\vM, \vM),\,(\vM, \vM),\,(\vM, \vM - 1) & 1 \\
B_{\mathrm{lift}} & \vM - 1 & (\vM - 1, 1),\,(\vM, 2\vM),\,(\vM, \vM - 1) & 1 \\
C_{\mathrm{lift}} & 1       & (1, \vM + 1),\,(\vM - 1, 2\vM - 1),\,(2\vM, \vM - 1) & 1 \\
F_1               & \vM + 1 & (\vM, \vM - 2),\,(\vM - 1, 1),\,(\vM, 2\vM - 1) & 1 \\
G_1               & \vM - 1 & (\vM - 1, 1),\,(\vM + 1, \vM - 1),\,(2\vM - 1, \vM) & -1
\end{array}
\]
The four $j = 1$ families satisfy
$\aK^{(3)} - 2^{\vM}\aI^{(3)} = -b$, while $G_1$ ($j = -1$) satisfies the
rational synchronisation $\aK^{(3)} - 2^{-\vM}\aI^{(3)} = 2^{-\vM} b$. In
addition there is a $\vM$-dependent number of further $q$-specific level-$4$
sub-classes --- six at $\vM = 3$, four at $\vM = 4$ --- which are not universal
in $\vM$; their count is not uniform in $\vM$.
\end{theorem}

\begin{proof}
For each of the five displayed exponent sequences, running the recursion of
Theorem~\ref{thm:recursion} symbolically in $\vM$ from $(\lambda_0, D_0) =
(2^{v_0}, -1)$ produces a pre-synchronisation state $(\lambda_3, D_3) =
(2^{j\vM}, -(2^{j\vM} - 1)/q)$ with the listed $j$; for $j = 1$ this is
$(2^{\vM}, -1)$, giving $\aK^{(3)} - 2^{\vM}\aI^{(3)} = -b$, and for $G_1$
the value $j = -1$ gives $\lambda_3 = 2^{-\vM}$ and $D_3 = -(2^{-\vM} - 1)/q
= 2^{-\vM}$, the displayed rational synchronisation --- the first sub-class
requiring a negative signed index, which is why the index must be taken in
$\Z$ rather than $\N$. The $q$-specific families are found by the same
recursion for fixed small $\vM$ but do not admit a closed form uniform in
$\vM$, and their number varies with $\vM$ (six at $\vM = 3$, four at
$\vM = 4$). The exact count of $q$-specific families is not used in any later
bound --- the density results of Section~\ref{sec:density} only require the
listed families to be present and pairwise disjoint, so the lower bounds hold
regardless of how many further families exist. The symbolic sync identities of
the five universal families are verified by \srcpath{code/python/recursion.py}; the
exhaustive level-$4$ sub-class count --- $11$ at $\vM = 3$ and $9$ at
$\vM = 4$, whence six and four $q$-specific families respectively --- is
verified by \srcpath{code/python/stage4\_enum.py}.
\end{proof}

% =============================================================================
\section{The node-lift theorem}\label{sec:nodelift}
% =============================================================================

\begin{theorem}[Node lift]\label{thm:node-lift}
Let $q = 2^{\vM} - 1$. The recursion of Theorem~\ref{thm:recursion} has a
unique fixed point carrying the synchronisation constant $-b$ (i.e.\ with
$D = -1$),
\[
(\lambda^*, D^*) = (2^{\vM + 1}, -1),
\]
attained by the step $(\vk, \vi) = (\vM, \vM)$. If a level-$k$ sub-class
passes through this $A$-node at some position $\ell$, then inserting one
further fixed-point step $(\vM, \vM)$ at position $\ell$ produces a level-$(k+1)$
sub-class with the same synchronisation constant $-b$. Consequently there are
four iterated families
\[
A_k\ (k \ge 2), \quad B_k\ (k \ge 4), \quad C_k\ (k \ge 5), \quad F_{1,k}\ (k \ge 4),
\]
each satisfying $\aK^{(k-1)} - 2^{\vM}\aI^{(k-1)} = -b$, generated by
iterating the $A$-node from the universal families of
Sections~\ref{sec:stage3}--\ref{sec:stage4} that thread it.
\end{theorem}

\begin{proof}
At the state $(\lambda, D) = (2^{\vM + 1}, -1)$ a step with
$(\vk, \vi) = (\vM, \vM)$ has $\vi - \vk = 0$, so $\lambda' = \lambda$, and
\[
E = 1 - \lambda - q D = 1 - 2^{\vM + 1} + q
  = 1 - 2^{\vM + 1} + (2^{\vM} - 1) = -2^{\vM},
\]
whence $D' = E/2^{\vk} = -2^{\vM}/2^{\vM} = -1$. Thus
$(\lambda', D') = (2^{\vM + 1}, -1) = (\lambda, D)$: the $A$-node is fixed.
It is moreover the unique fixed point carrying the synchronisation constant
$-b$, i.e.\ with $D = -1$. Indeed, solving $\lambda' = \lambda$, $D' = D$ for
such a state: $\lambda' = \lambda$ forces $\vi = \vk$ (as $\lambda' =
\lambda 2^{\vi - \vk}$ with $\lambda \ne 0$), and then $D' = D = -1$ reads
$E = 1 - \lambda + q = -2^{\vk}$, i.e.\
\[
\lambda = 1 + q + 2^{\vk} = 2^{\vM} + 2^{\vk}
\]
(using $1 + q = 2^{\vM}$). A sum of two powers of two is itself a power of
two if and only if the two exponents are equal, so $\vk = \vM$ and
$\lambda = 2^{\vM + 1}$; the $A$-node is the only such state, with no
restriction on the valuation range required.

Because the inserted step neither changes $(\lambda, D)$ nor the eventual
synchronising value $(2^{\vM}, -1)$ reached one step later, a level-$k$
sub-class threading the node at position $\ell$ lifts to a level-$(k+1)$
sub-class with synchronisation constant $-b$ unchanged. Four universal
families thread the $A$-node: $A_2$ and $F_{1,4}$ start \emph{at} it (their
$v_0 = \vM + 1$ gives $(\lambda_0, D_0) = (2^{\vM + 1}, -1)$), while $B_4$ and
$C_5$ reach it after some steps. Iterating the fixed-point insertion from each
generates the towers $A_k\ (k \ge 2)$, $F_{1,k}\ (k \ge 4)$, $B_k\ (k \ge 4)$,
$C_k\ (k \ge 5)$ at the stated minimal levels, all with synchronisation
constant $-b$. We do not claim these are the only node-threading sub-classes
at every level: the lower bounds of Section~\ref{sec:density} use only their
presence and pairwise disjointness, and the asymptotic upper bound is
conditional on the completeness hypothesis of Theorem~\ref{thm:asymptotics}.
The fixed-point computation and the iterated synchronisation identities are
verified symbolically in $\vM$ up to $k = 8$ by \srcpath{code/python/recursion.py}.
\end{proof}

\begin{remark}
The level-$4$ family $C_{\mathrm{lift}}$ of Theorem~\ref{thm:stage4} (valuation
sequence ending $(\vk^{(2)}, \vi^{(2)}) = (2\vM, \vM - 1)$) and the base $C_5$
of the C-tower (ending $(2\vM, \vM)$, synchronising at level $5$) are distinct
families that happen to share the letter $C$; only $C_5$ threads the $A$-node.
\end{remark}

% =============================================================================
\section{Density lower bound}\label{sec:density}
% =============================================================================

\begin{remark}[2-adic reading of the density]\label{rem:haar}
By the lift theorem of \cite{ObstructionResiduesFamily2026} the obstruction
sets are lift-stable, so they define an open subset of the $2$-adic integers
$\Ztwo$ (an increasing union of residue cylinders, the density being a
non-decreasing limit), and $\cWq$ is its normalised Haar measure --- the negative-Mersenne
instance of the $2$-adic density placed in the Bernstein--Lagarias tradition
\cite{BernsteinLagarias1996} (for the $2$-adic background see
\cite{Robert2000}). The bounds below are therefore measure lower bounds on a
fixed subset of $\Ztwo$.
\end{remark}

\begin{theorem}[Density lower bound]\label{thm:atom-density}
Let $q = 2^{\vM} - 1$ with $\vM \ge 2$. Then
\[
\cWq \ge \frac{1}{2q}.
\]
\end{theorem}

\begin{proof}
By Corollary~\ref{cor:mersenne-sync} the principal Mersenne-synchronising
class is exactly the residues with $\vtwo(r + b) = \vM$ (those synchronise in a
single step), a cylinder of density $2^{-\vM} - 2^{-\vM - 1} = 2^{-\vM - 1}$,
and by
Theorem~\ref{thm:node-lift} the iterated family $A_k$ ($k \ge 2$) occupies,
for each $k$, a residue cylinder of density $2^{-k\vM - 1}$; identifying the
principal class as the $k = 1$ term of the same geometric pattern, these
cylinders are pairwise disjoint, because the synchronisation level is by
definition the \emph{least} index at which the two tracks agree and is thus a
well-defined function of $r$, so each obstruction lies in the cylinder of
exactly one $k$. Summing,
\[
\cWq \ge \sum_{k = 1}^{\infty} 2^{-k\vM - 1}
      = \frac{2^{-\vM - 1}}{1 - 2^{-\vM}}
      = \frac{2^{-\vM - 1}\, 2^{\vM}}{2^{\vM} - 1}
      = \frac{1}{2(2^{\vM} - 1)}
      = \frac{1}{2q}. \qedhere
\]
\end{proof}

\begin{corollary}[Refined lower bound]\label{cor:atom-density-extended}
Let $\vM \ge 3$. Counting, beyond the principal class and the $A$-family,
the full $B$-, $C$-
and $F$-families (each the geometric tower of its members across
synchronisation levels) and the family $G_1$,
\[
\cWq \ge \frac{1}{2q}
       + \frac{2^{-2\vM}}{q}
       + \frac{2^{-2\vM - 1}}{q}
       + \frac{2^{-3\vM}}{q}
       + 2^{-4\vM},
\]
the four correction terms being the totals of the $B$-, $C$-, $F$- and
$G_1$-families respectively.
\end{corollary}

\begin{proof}
Each universal family is a tower of residue cylinders, one per synchronisation
level $k$ at and above its minimal level, and each cylinder's density is read
off from its modulus $2^{\max(\Vk^{(k-1)},\, v_0 + \Vi^{(k-1)}) + 1}$. By
synchronisation level the towers are pairwise disjoint, and within a level
distinct families have distinct valuation data (the nesting $v_0$ \emph{together
with} the valuation sequence --- $v_0$ alone does not separate them, e.g.\
$A_{\mathrm{lift}}$ and $F_1$ share $v_0 = \vM + 1$), so all the cylinders below
are disjoint. The members of each family letter have a fixed per-level density,
giving a geometric series with ratio $2^{-\vM}$ and
$\sum_{k \ge k_0} 2^{-k\vM} = 2^{-k_0\vM}\,2^{\vM}/q$:
\[
\renewcommand{\arraystretch}{1.3}
\begin{array}{l|l|l|l}
\text{family} & \text{min.\ level } k_0 & \text{per-level density} & \text{total} \\ \hline
B & 3 & 2^{-k\vM} & \sum_{k\ge3} 2^{-k\vM} = 2^{-2\vM}/q \\
C & 3 & 2^{-(k\vM + 1)} & \sum_{k\ge3} 2^{-(k\vM+1)} = 2^{-2\vM - 1}/q \\
F & 4 & 2^{-k\vM} & \sum_{k\ge4} 2^{-k\vM} = 2^{-3\vM}/q \\
G_1 & 4 & 2^{-4\vM} & 2^{-4\vM}
\end{array}
\]
(the $A$-family, levels $k \ge 1$ with per-level density $2^{-(k\vM+1)}$, already
sums to $1/(2q)$ in Theorem~\ref{thm:atom-density}). Adding the four totals to
$1/(2q)$ gives the stated bound. The level-$3$ members of the $B$- and
$C$-families are the families $B$, $C$ of Theorem~\ref{thm:stage3}; the
level-$4$ members include $B_{\mathrm{lift}}$, $C_{\mathrm{lift}}$, $F_1$, $G_1$
of Theorem~\ref{thm:stage4}; the higher members are the node-lift towers of
Theorem~\ref{thm:node-lift}. The per-family densities and the identity of the
total with the displayed bound are checked by
\srcpath{code/python/refined\_decomposition.py}, and the table values by
\srcpath{code/python/atom\_density.py}. (At $\vM = 2$ the tower decomposition
degenerates --- families $B$ and $C$ then share $v_0 = 1$ --- so the bound is
stated for $\vM \ge 3$; the $\vM = 2$ entry of Table~\ref{tab:density} is the
formula evaluated at $\vM = 2$, which remains a lower bound on $\cWq$ because
the monotone finite density of Theorem~\ref{thm:atom-density} already exceeds
it.)
\end{proof}

\begin{table}[t]
\centering
\renewcommand{\arraystretch}{1.15}
\begin{tabular}{ccccc}
\toprule
$\vM$ & $q$ & $1/(2q)$ & refined bound & $\lvert\Obsq\rvert/2^L$, $L = 16$ \\
\midrule
$2$ & $3$  & $0.1667$ & $0.2070$ & $0.2498$ \\
$3$ & $7$  & $0.0714$ & $0.0753$ & $0.0770$ \\
$4$ & $15$ & $0.0333$ & $0.0338$ & $0.0338$ \\
$5$ & $31$ & $0.0161$ & $0.0162$ & $0.0162$ \\
\bottomrule
\end{tabular}
\caption{The simple bound $1/(2q)$, the refined bound of
Corollary~\ref{cor:atom-density-extended}, and the finite density
$\lvert\Obsq\rvert/2^L$ at $L = 16$ (with $b = 1$). Both bounds are lower
bounds on the limit $\cWq$; the finite density is itself a monotone lower
bound (Theorem~\ref{thm:atom-density}), still increasing for small $\vM$ and
essentially converged for $\vM \ge 4$, where it agrees with the refined bound
to the displayed precision. (For large $\vM$ the finite density at a fixed
level can sit just below the refined bound, since both approach $\cWq$ from
below.)}
\label{tab:density}
\end{table}

Table~\ref{tab:density} compares both bounds with the finite density; the
values are recomputed and checked against the table by
\srcpath{code/python/atom\_density.py}.

% =============================================================================
\section{The observability bound}\label{sec:jbound}
% =============================================================================

\begin{theorem}[Observability bound]\label{thm:j-bound}
Let $r$ be a residue in a level-$k$ sub-class with signed synchronisation
index $j$. Then $r$ has modulus at least
$2^{\,k + \lvert j \rvert\, \vM}$, so it contributes to $\Obsq$ only for
\[
L \ge k + \lvert j \rvert\, \vM, \qquad \text{equivalently} \qquad
\lvert j \rvert \le \frac{L - k}{\vM} .
\]
\end{theorem}

\begin{proof}
Write $A := \Vk^{(k-1)}$ and $B := v_0 + \Vi^{(k-1)}$ for the cumulative
modulus consumption of the two tracks up to synchronisation. The signed index
satisfies $j\vM = B - A$, so $\lvert A - B \rvert = \lvert j \rvert\, \vM$;
since synchronisation at level $k$ takes $k - 1$ steps and each consumes at
least one unit of valuation on each track --- the iterates $\aK^{(i)},
\aI^{(i)}$ are odd and $b$ is odd, so $-q\,\aK^{(i)} + b$ and
$-q\,\aI^{(i)} + b$ are even, forcing $\vk^{(i)}, \vi^{(i)} \ge 1$ --- both
$A$ and $B$ are sums of at
least $k - 1$ positive terms, so $\min(A, B) \ge k - 1$, hence
\[
\max(A, B) = \min(A, B) + \lvert A - B \rvert \ge (k - 1) + \lvert j \rvert\, \vM .
\]
The residue carrying this sub-class has modulus $2^{\max(A,B) + 1}$, so it is
visible at level $L$ only when $L \ge \max(A, B) + 1 \ge k + \lvert j \rvert\,
\vM$, which is the claim. The bound is exercised against the enumerated
sub-classes by \srcpath{code/python/j\_bound.py}.
\end{proof}

% =============================================================================
\section{Asymptotic sharpness}\label{sec:asymptotics}
% =============================================================================

\begin{theorem}[Asymptotics of the density]\label{thm:asymptotics}
Let $q = 2^{\vM} - 1$. The refined bound of
Corollary~\ref{cor:atom-density-extended} gives the lower bound
\[
\cWq \ge \frac{1}{2q} + \frac{3}{2}\,\frac{2^{-2\vM}}{q}
\]
(its further terms, of order $2^{-3\vM}/q$, being positive).
Suppose moreover the \emph{completeness hypothesis} that, as
$\vM \to \infty$, the universal families of
Sections~\ref{sec:stage3}--\ref{sec:nodelift} exhaust the obstruction density
up to order $2^{-3\vM}/q$:
\[
\cWq - (\text{universal-family total})
   = O\!\left(\frac{2^{-3\vM}}{q}\right),
\]
equivalently, the aggregate contribution of the $q$-specific and
level-$(\ge 5)$ families is $O(2^{-3\vM}/q)$ (related to the sub-class growth
of Conjecture~\ref{conj:subclass-count}). Then the
matching upper bound holds and
\[
\cWq = \frac{1}{2q} + \frac{3}{2}\,\frac{2^{-2\vM}}{q}
     + O\!\left(\frac{2^{-3\vM}}{q}\right)
     = \frac{1}{2q}\big(1 + O(2^{-\vM})\big),
\]
so the simple lower bound $1/(2q)$ of Theorem~\ref{thm:atom-density} is
asymptotically sharp. The hypothesis is corroborated by the density table
(Table~\ref{tab:density}) but is not proved here; the sub-class growth it
relies on is open (Conjecture~\ref{conj:subclass-count}).
\end{theorem}

\begin{proof}
The lower bound is Corollary~\ref{cor:atom-density-extended}: the principal
class and the $A$-family sum exactly to $1/(2q)$
(Theorem~\ref{thm:atom-density}), and the next-order families contribute its
four correction terms --- the $B$-family gives $2^{-2\vM}/q$, the $C$-family
$2^{-2\vM - 1}/q$, the $F$-family $2^{-3\vM}/q$, and $G_1$ gives $2^{-4\vM}$.
The two leading ones (the $B$- and $C$-families) sum to
\[
\frac{2^{-2\vM}}{q} + \frac{2^{-2\vM - 1}}{q}
  = \Big(1 + \tfrac12\Big)\frac{2^{-2\vM}}{q}
  = \frac{3}{2}\,\frac{2^{-2\vM}}{q},
\]
while the $F$-family ($2^{-3\vM}/q$) and $G_1$ ($2^{-4\vM} =
\Theta(2^{-3\vM}/q)$, since $q = 2^{\vM} - 1$) are of the next order. For the matching upper bound, the
order-$2^{-2\vM}/q$ mass is carried by the $B$- and $C$-families (their full
towers), each level-$4$ $q$-specific family contributing only at the next order
$2^{-4\vM} = O(2^{-3\vM}/q)$ and their number being bounded
(Conjecture~\ref{conj:qspecific}). Under the stated hypothesis the level-$(\ge 5)$ and any
remaining $q$-specific families add at most $O(2^{-3\vM}/q)$ in aggregate, so
no further term arises at order $2^{-2\vM}/q$ and the upper bound matches the
lower one; the relative correction to $1/(2q)$ is then $O(2^{-\vM})$. We stress
that this aggregate-tail hypothesis is not established: the number of level-$k$
sub-classes grows (Conjecture~\ref{conj:subclass-count}), and we have no
a-priori bound forcing the tail below $O(2^{-3\vM}/q)$; the conclusion is
therefore conditional. This is an asymptotic ($\vM \to \infty$) statement, not
a finite enumeration; the finite lower bound it sharpens, and the table values
feeding the constant $\tfrac32$, are checked by
\srcpath{code/python/atom\_density.py}.
\end{proof}

% =============================================================================
\section{Open problems}\label{sec:open}
% =============================================================================

\begin{conjecture}[No uniform $q$-specific form]\label{conj:qspecific}
At each $\vM$ the level-$k$ hierarchy contains finitely many $q$-specific
sub-classes (e.g.\ four at $\vM = 4$, level $4$) beyond the universal
families, but there is no closed form for them uniform in $\vM$.
\end{conjecture}

\begin{conjecture}[Sub-class growth]\label{conj:subclass-count}
The number $\#\Sigma_k(\vM)$ of level-$k$ sub-classes grows exponentially in
$k$ with ratio between $3$ and $4$; no second-order linear recurrence fits
beyond three data points.
\end{conjecture}

\begin{conjecture}[Spectral exponent]\label{conj:alpha}
The growth exponent $\alpha(\vM)$ governing $\lvert \Atomq \rvert$, the
spectral radius of the transfer matrix of the recursion, has no closed form;
the empirical estimates remain $L$-dependent within the tested range. An
operator-theoretic route in the spirit of the Wirsching operator for the
$3n+1$ map \cite{Wirsching1998} is the natural setting for this question.
\end{conjecture}

\begin{conjecture}[Non-Mersenne density]\label{conj:nonmersenne}
For non-Mersenne $q$ the principal class is absent
(Corollary~\ref{cor:nonmersenne}) and, empirically,
$\cWq = \Theta\big(4^{-\ord_q(2)}\big)$, matching the generic positive-multiplier
scale of \cite{ObstructionResiduesFamily2026} rather than the
$\Theta(2^{-\vM})$ scale of the Mersenne ladder.
\end{conjecture}

\begin{conjecture}[Sign conjugation]\label{conj:conjugation}
The positive-multiplier obstructions are not contained in the
negative-multiplier ones, $\Obs^{(+q)} \not\subseteq \Obs^{(-q)}$; whether an
arithmetic bijection (rather than the failed identity inclusion) relates the
obstructions of the two signs is open.
\end{conjecture}

% =============================================================================
\section*{Funding and competing interests}
% =============================================================================

The author received no funding for this research and declares no competing
interests.

\section*{Use of AI tools}

This manuscript was prepared with substantial assistance from generative AI
tools. Anthropic's Claude assisted with developing and formulating the proofs,
drafting and refining the exposition, and writing the accompanying verification
code; OpenAI's ChatGPT was used additionally for independent second-opinion
reviews of the arguments. All definitions, theorems, proofs, and results were
reviewed and independently checked by the author and are backed by the
reproducible verification and falsification code accompanying this manuscript. The author takes
full and sole responsibility for the correctness and content of this
manuscript. No AI system is listed as an author.

% =============================================================================
\bibliographystyle{amsplain}
\bibliography{references}

\end{document}
